% !PRINTED 564
\weberpage{564}
%<<CONT
hence $\xi$, $\xi_2$, $\xi_4$ are the roots of the cubic equation
\weq{16}{x^3 - \zeta x^2 - x - 1 + \zeta = 0.}

\begin{center}
\S.\,169.

\gesperrt{Reduction of the cyclotomic equation to}\\
\gesperrt{pure equations.}
\end{center}

\gesperrt{Gauss}, already in his first exposition in the disq.
arithm., reduced the cyclotomic equations directly to \gesperrt{pure equations}
by the use of the resolvents\footnote{\gesperrt{Gauss}, disq. arithm. art. 359, 360; disq. circa aequationes puras ulterior evolutio, Werke Bd. II. \gesperrt{Lagrange}, r\'es. des \'equations num\'eriques. \gesperrt{Jacobi}, ``Ueber die Kreistheilung und ihre Anwendung in der Zahlentheorie''. Werke Bd. 6. \gesperrt{Kummer}, Crelle's Journal, Bd. 35 and Abhandl. d. Berl. Akademie 1856. \gesperrt{Eisenstein} and \gesperrt{Cauchy} are also to be mentioned. The theory of the division of the circle is presented in connected form, accompanied by the historical references, in the book by \gesperrt{Bachmann}, ``Die Lehre von der Kreistheilung''. Leipzig 1872.}. To
carry this out, we set
\[ n - 1 = m, \]
and denote, as in the preceding section, by $\varepsilon$ a primitive
$m$th root of unity, and by $r, r_1, r_2 \ldots r_{m-1}$ the $n$th roots of unity
 in the order fixed above. The \gesperrt{Lagrange}
resolvents are then
\[ (\varepsilon^\lambda, r) = \sum_{0,\,n-2}^{h} \varepsilon^{\lambda h} r_h = \sum_{0,\,n-2}^{h} \varepsilon^{\lambda h} r^{g^h}, \]
in which $\lambda$ is an arbitrary exponent taken with respect to the modulus $m$.
If we set $g^h = s$, so that $h \equiv \mathrm{ind}\,s \pmod{m}$, we may
also write
\weq{1}{(\varepsilon^\lambda, r) = \sum_{1,\,n-1}^{s} \varepsilon^{\lambda\,\mathrm{ind}\,s} r^s.}

The basis for the further reduction of these expressions
is the computation of the product of two such resolvents.

We denote by $\mu$ a second number like $\lambda$, and understand
by $s$ and $t$ two symbols which, independently of one another, each
run through a system of residues with respect to the modulus $n$, with the exclusion
of zero.
Then
\[ (\varepsilon^\lambda, r)(\varepsilon^\mu, r) = \sum^{s} \sum^{t} r^{(s+t)} \varepsilon^{\lambda\,\mathrm{ind}\,s + \mu\,\mathrm{ind}\,t}. \]
